\documentclass[10pt,a4paper]{amsart}
\usepackage[margin=19mm]{geometry}
\usepackage{amsmath,amssymb,mathtools}
\usepackage[scr=rsfs,cal=euler]{mathalfa}
\usepackage{xfrac}
\usepackage[unicode,hidelinks,bookmarks=false]{hyperref}
\newcommand{\ie}{i.e.\ }
\newcommand{\cf}{cf.\ }
\newcommand{\resp}{respectively\ }
\newcommand{\Subsection}{\subsection}
\providecommand{\nobreakdash}{\nobreak\hbox{-}}
\setlength{\emergencystretch}{2em}
\multlinegap0pt
\usepackage{bm}
\usepackage{bbm}
\usepackage{dsfont}
\usepackage{xspace}
\usepackage{cancel}
\usepackage{mathtools}
\usepackage{amsthm}
\makeatletter
\@ifundefined{Theorem}{\newtheorem{Theorem}{Theorem}}{}
\@ifundefined{Definition}{\newtheorem{Definition}[Theorem]{Definition}}{}
\@ifundefined{Proposition}{\newtheorem{Proposition}[Theorem]{Proposition}}{}
\makeatother
\newcommand{\Pf}{\operatorname{Pf}}
\newcommand{\R}{{\mathbb R}}
\newcommand{\cA}{\mathcal{A}}
\newcommand{\cP}{{\mathcal P}}
\newcommand{\cS}{\mathcal{S}}
\newcommand{\dd}{\Theta}
\newcommand{\ind}{{\bf{1}}}
\newcommand{\lst}[1]{[\![#1 ]\!]}
\newcommand{\sgn}{\operatorname{sgn}}
\title[Solvable Families of Random Block Tridiagonal Matrices]{Solvable Families of Random Block Tridiagonal Matrices}
\author{Brian Rider, Benedek Valk\'o}
\date{Source: arXiv:2412.04579v3 (CC BY 4.0)}
\begin{document}
\maketitle

\noindent\emph{Source and licence.} Solvable Families of Random Block Tridiagonal Matrices. Version arXiv:2412.04579v3 (CC BY 4.0), distributed under the Creative Commons Attribution 4.0 International licence, as recorded in the arXiv licence metadata for this exact version and shown on the abstract page https://arxiv.org/abs/2412.04579v3. Full rights notice: licenses/arxiv-2412.04579-CC-BY-4.0.txt.\par\smallskip

\noindent\textit{Editorial scope.} Counted unit: the Proposition whose source label is \texttt{prop:det2+Pf} in the article (source0.tex lines 1380--1386, source label \texttt{prop:det2+Pf}), delivered together with the article's own complete original proof of it (source0.tex lines 1389--1440, one \texttt{proof} environment of 52 lines). No mathematics is written, repaired or re-derived: statement and proof are the article's own text line for line. Required context retained uncounted: shorthand (lines 917--928); eq:defdelta (lines 1229--1231); eq:Pfaffiandef (lines 1366--1369). Every definition, display and label of those lines is retained unchanged. Omitted from this package and not redistributed: 1--916, 929--1228, 1232--1365, 1370--1379, 1387--1388, 1441--2093. Nothing inside a retained excerpt is omitted or altered: every retained body is delivered exactly as the article's lines are written, so no omission receipt of any kind accompanies this package. Citations: the retained excerpts carry no citation command at all, so no bibliography entry of the article is transcribed and no reference is redistributed. Reference adaptation: none was needed and none was made; every reference inside the delivered unit resolves inside it, so no unresolved reference or citation is produced. Printed slips kept literal: no printed word, symbol, hypothesis or number of the counted statement or of its proof was altered. Layout and class: the article's own class is \texttt{sigma} with packages including ifpdf; the wrapper is the lane's pinned \texttt{amsart} editorial wrapper at 10pt on A4, which restates the theorem-like environments the retained text needs and adds the article's own macro definitions that the retained text needs (\texttt{\textbackslash Pf}, \texttt{\textbackslash R}, \texttt{\textbackslash cA}, \texttt{\textbackslash cP}, \texttt{\textbackslash cS}, \texttt{\textbackslash dd}, \texttt{\textbackslash ind}, \texttt{\textbackslash lst}, \texttt{\textbackslash sgn}), with extra wrapper packages (bm, bbm, dsfont, xspace, cancel, mathtools). The delivered numbering is the unit's own. Assets: no retained span contains a figure environment, a graphics-inclusion command, a PDF-page inclusion, a table or a TikZ picture; no image file is copied into this package, the package declares no asset dependency and no image file is redistributed with it. Licence: version arXiv:2412.04579v3 of this article is distributed under the Creative Commons Attribution 4.0 International licence, as recorded in the arXiv licence metadata for this exact version and shown on the abstract page https://arxiv.org/abs/2412.04579v3. A full rights notice is delivered at licenses/arxiv-2412.04579-CC-BY-4.0.txt. Characters: every retained excerpt is pure ASCII, so the delivered engine, XeLaTeX with its default Latin Modern text font, reports no missing character.
\begin{Definition}
 For $\underline{\cA}=(\cA_1,\dots,\cA_r)\in \cP_{r,n}$, let $\sigma_{\underline{\cA}}\in S_{rn}$
 denote the following permutation of the set $\lst{rn}$
$
\sigma_{\underline{\cA}}(a+(m-1)n)= $ the $a$-th largest element of $ \cA_m$, for $a\in \lst{n}$, $ m\in \lst{r}$.
Now when $r=2$, the elements of $(\cA_1, \cA_2)$ of $\cP_{2,n}$ can be identified just with the subsets $\cA\subset \lst{2n}$ of cardinality~$n$. Using $\cA'$ for the complement of $\cA$ in $\lst{2n}$ throughout, we further denote
\begin{equation}
\label{shorthand}
\sigma_{{\cA}} := \sigma_{\cA, \cA'} \in S_{2n}.
\end{equation}

\end{Definition}

\begin{align}\label{eq:defdelta}
 \dd(\cS_1,\cS_2):= \prod_{s_1\in \cS_1} \prod_{s_2\in \cS_2}(\lambda_{s_1}-\lambda_{s_2}),
\end{align}

\begin{align}\label{eq:Pfaffiandef}
 \Pf(M)=\frac{1}{2^n n!} \sum_{\sigma\in S_{2n}} \sgn(\sigma) \prod_{i=1}^n M_{{\sigma(2i-1)},{\sigma(2i)}}
 = \sum_{\alpha\in \Pi_{2n}} \sgn(\pi_\alpha) \prod_{i=1}^n M_{\alpha_{i,1},\alpha_{i,2}},
\end{align}

\begin{Proposition}\label{prop:det2+Pf}
Let ${\lambda}\in \R^{2n}$. Then we have
 \begin{align*}%\label{eq:det2+Pf}
 \sum_{|\cA|=n,\, \cA\subset\lst{2n}} \Delta(\cA)^2 \Delta (\cA')^2&=(-2)^n \Delta (\lambda) \Pf\left(\frac{\ind_{i\neq j}}{\lambda_i-\lambda_j}\right).
 \end{align*}
The right side is defined as the appropriate limit if the $\lambda_i$'s are not all distinct.
\end{Proposition}

\begin{proof} Assuming that the $\lambda_i$'s are distinct, the definition \eqref{eq:Pfaffiandef} gives
\begin{align*}
 \Pf\left(\frac{\ind_{i\neq j}}{\lambda_i-\lambda_j}\right)&
 %=\frac{1}{2^n n!} \sum_{\sigma\in S_{2n}} \sgn(\sigma) \prod_{i=1}^n
 %\frac{1}{\lambda_{\sigma(2i-1)}-\lambda_{\sigma(2i)}}
 =\sum_{\alpha\in \Pi_{2n}} \sgn(\pi_\alpha) \prod_{i=1}^n \frac{1}{\lambda_{\alpha_{i,1}}-\lambda_{\alpha_{i,2}}}.
\end{align*}
A matching $\alpha \in \Pi_{2n}$ can be encoded with a pair $(\cA, \tau)$ where $|\cA|=n$, $\cA\subset \lst{2n}$ and $\tau\in S_n$. Indeed, $\cA=\{\alpha_{1,1}, \dots, \alpha_{n,1}\}$ and the permutation $\tau\in S_n$ that takes the ordered version of $\cA'$ into $(\alpha_{1,2}, \dots, \alpha_{n,2})$ identifies the matching $\alpha$. Note that $\sgn(\pi_\alpha)$ is equal to $\sgn(\sigma_{\cA} )\sgn(\eta_n)$ where $\sigma_{\cA}$ is defined in \eqref{shorthand}, and $\eta_n$ is the permutation $(1,n+1,2, n+2, \dots, n, 2n)$. This leads~to
\begin{align}
 \label{eq:Pf4_2}
\Pf\left(\frac{\ind_{i\neq j}}{\lambda_i-\lambda_j}\right)=\sgn(\eta_n) \sum_{|\cA|=n,\, \cA\subset \lst{2n}} \sgn(\sigma_{\cA}) \sum_{\tau\in S_n} \sgn(\tau) \prod_{i=1}^n \frac{1}{\lambda_{a'_{\tau(i)}}-\lambda_{a_i}},
\end{align}
where we used $a_i$, $a_i'$, $i\in \lst{n}$, to denote the (ordered) elements of $\cA$ and $ \cA'$, respectively.

For a given $|\cA|=n$, $\cA\subset \lst{2n}$, we have
\[
\Delta(\lambda)=\Delta(\cA)\Delta(\cA') \prod_{\substack{1\le i<j\le 2n\\ \left|\{i,j\}\cap \cA\right|=1}} (\lambda_j-\lambda_i).
\]
By \eqref{eq:defdelta}, we have
\[
\prod_{\substack{1\le i<j\le 2n\\ \left|\{i,j\}\cap \cA\right|=1}} (\lambda_j-\lambda_i)=(-1)^{k_\cA}
\dd(\cA, \cA')
\]
 with $k_{\cA}$ the number of pairs $(i,j)$ with $1\le i<j\le 2n$ and $i\in \cA$, $j\in \cA'$. Observe then that~${n^2-k_{\cA}}$ is exactly the number of inversions of the permutation $\sigma_{\cA}$.
Hence,
 \begin{align}
 \label{eq:det_exp_A_Ac}
\Delta
%_{{\lambda}}
(\lambda)=(-1)^n \sgn(\sigma_{\cA}) \Delta
%_{{\lambda}}
(\cA)\Delta
%_{{\lambda}}
(\cA') \dd(\cA, \cA').
\end{align}
Using this identity, we get
\begin{align}
 \sum_{|\cA|=n,\, \cA\subset\lst{2n}} \frac{\Delta
 %_{{\lambda}}
 (\cA)^2 \Delta
 %_{{\lambda}}
 (\cA')^2}{ \Delta(\lambda)}&=(-1)^n \sum_{|\cA|=n,\, \cA\subset\lst{2n}}
 \sgn(\sigma_{\cA})
 \frac{\Delta(\cA) \Delta(\cA')}{\dd(\cA, \cA')}\nonumber \\
 &=(-1)^{n+\binom{n}{2}}\sum_{|\cA|=n,\, \cA\subset\lst{2n}} \sgn(\sigma_{\cA}) C_{\cA, \cA'}
\nonumber
 \\
 &=(-1)^{n+\binom{n}{2}}\sum_{|\cA|=n,\, \cA\subset\lst{2n}} \sgn(\sigma_{\cA}) \sum_{\tau\in S_n} \sgn(\tau) \frac{1}{\lambda_{a_i}-\lambda_{a'_{\tau(i)}}}, \label{eq:det4_123}
\end{align}
where in the last line we expanded the Cauchy determinant, and used that $a_i$, $a_i'$, $i\in \lst{n}$, for the ordered elements of $\cA$, $\cA'$, respectively.
Note that \smash{$\sgn(\eta_n)=(-1)^{\binom{n}{2}}$}, hence by comparing~\eqref{eq:det4_123} with~\eqref{eq:Pf4_2} the statement follows.
\end{proof}

\end{document}
